% !TEX program = xelatex
% SY02-15 相对朗兰兹对偶（LAWRGE 2024 讲义）—— 主控导言区
% 结构照 SY01-29；定理与公式共号照原文（原文每节一个计数器）
\usepackage[a4paper,margin=2.6cm,top=2.8cm,bottom=2.8cm]{geometry}
% 原文公式编号排在左边（leqno）
\PassOptionsToPackage{leqno}{amsmath}
\usepackage{amsmath,amssymb,amsfonts,amsthm,mathtools}
\usepackage{mathrsfs}
\usepackage{bm}
\usepackage{graphicx}
\usepackage{booktabs}
\usepackage{array}
\usepackage{longtable}
\usepackage{enumitem}
\setlist{topsep=0.35em,itemsep=0.12em,parsep=0pt,partopsep=0pt}
\usepackage{xcolor}
\usepackage{tcolorbox}
\usepackage{hyperref}
\usepackage{xspace}
\usepackage{tikz}
\usepackage{epigraph}

\graphicspath{{figures/}}

% ---------- 字体与行距 ----------
\linespread{1.22}
\setlength{\parskip}{0.25em}
\emergencystretch=1.5em
\sloppy

% 中英之间已手工加空格，关掉 xeCJK 的自动胶，避免与全角标点叠加出空隙
\xeCJKsetup{CJKecglue={\hskip 0.24em plus 0.05em minus 0.05em}}

% ---------- 章节格式：原文作「1. 标题」「1.1. 标题」 ----------
\ctexset{
  section = {
    format = \Large\bfseries\raggedright,
    name = {,},
    number = \arabic{section},
    aftername = {.\enspace},
    beforeskip = 2.4ex plus .4ex minus .2ex,
    afterskip = 1.4ex plus .2ex,
  },
  subsection = {
    format = \large\bfseries\raggedright,
    name = {,},
    number = \arabic{section}.\arabic{subsection},
    aftername = {.\enspace},
    beforeskip = 1.8ex plus .3ex minus .2ex,
    afterskip = 1ex plus .2ex,
  },
  subsubsection = {
    format = \normalsize\bfseries\raggedright,
    name = {,},
    number = \arabic{section}.\arabic{subsection}.\arabic{subsubsection},
    aftername = {.\enspace},
  },
}

% ---------- 编号：公式与定理共用节内计数器；图按节 ----------
\numberwithin{equation}{section}
\numberwithin{figure}{section}
\setcounter{tocdepth}{2}
\setcounter{secnumdepth}{3}

% ---------- 超链接 ----------
\hypersetup{
  colorlinks=true,
  linkcolor=blue!45!black,
  citecolor=green!40!black,
  urlcolor=magenta!60!black,
  pdftitle={相对朗兰兹对偶（LAWRGE 2024 讲义）},
  pdfauthor={Ben-Zvi, Sakellaridis, Venkatesh（LAWRGE 2024 参会者笔记）},
}

% ---------- 定理类环境：一律 [equation]，与原文共号 ----------
\theoremstyle{plain}
\newtheorem{thm}[equation]{定理}
\newtheorem{lem}[equation]{引理}
\newtheorem{prop}[equation]{命题}
\newtheorem{cor}[equation]{推论}
\newtheorem{conj}[equation]{猜想}
\theoremstyle{definition}
\newtheorem{defn}[equation]{定义}
\newtheorem{exmp}[equation]{例}
\newtheorem{exer}[equation]{练习}
\newtheorem{claim}[equation]{断言}
\theoremstyle{remark}
\newtheorem{rem}[equation]{注}
\newtheorem*{conv}{约定}
\newtheorem*{notn}{记号}

% ---------- 译注 ----------
\newcounter{transnote}
\newenvironment{yizhu}
  {\par\medskip\refstepcounter{transnote}%
   \begin{tcolorbox}[colback=orange!4,colframe=orange!55!black,boxrule=0.5pt,
                     left=4pt,right=4pt,top=3pt,bottom=3pt,arc=1pt]
   \noindent\textbf{译注 \thetransnote：}}
  {\end{tcolorbox}\par\medskip}
% 译注脚注：不占用作者的脚注计数器（原文脚注 1,2,3… 必须与原文编号一致）
\newcommand{\zy}[1]{%
  \begingroup\renewcommand{\thefootnote}{}\footnotetext{\textcolor{orange!70!black}{【译注】}#1}\endgroup}

% ---------- 常用记号 ----------
\newcommand{\IC}{\mathbb{C}}
\newcommand{\IN}{\mathbb{N}}
\newcommand{\IZ}{\mathbb{Z}}
\newcommand{\IR}{\mathbb{R}}
\newcommand{\IQ}{\mathbb{Q}}
\newcommand{\IP}{\mathbb{P}}
\newcommand{\IA}{\mathbb{A}}
\newcommand{\IF}{\mathbb{F}}
\newcommand{\IG}{\mathbb{G}}
\newcommand{\hott}{\widehat{\otimes}}
\newcommand{\la}{\langle}
\newcommand{\ra}{\rangle}
\newcommand{\half}{\tfrac{1}{2}}
\newcommand{\abs}[1]{\lvert #1 \rvert}
\newcommand{\norm}[1]{\lVert #1 \rVert}
\newcommand{\eps}{\varepsilon}
\newcommand{\id}{\mathrm{id}}
\newcommand{\pt}{\mathrm{pt}}
\newcommand{\ev}{\mathrm{ev}}
\newcommand{\Fr}{\mathrm{Fr}}

\DeclareMathOperator{\Bun}{Bun}
\DeclareMathOperator{\Loc}{Loc}
\DeclareMathOperator{\Shv}{Shv}
\newcommand{\QC}{\mathcal{QC}}
\DeclareMathOperator{\QCoh}{QCoh}
\DeclareMathOperator{\Coh}{Coh}
\DeclareMathOperator{\IndCoh}{IndCoh}
\DeclareMathOperator{\Perf}{Perf}
\DeclareMathOperator{\Rep}{Rep}
\DeclareMathOperator{\Hom}{Hom}
\DeclareMathOperator{\End}{End}
\DeclareMathOperator{\Aut}{Aut}
\DeclareMathOperator{\Sym}{Sym}
\DeclareMathOperator{\Func}{Fun}
\newcommand{\gquot}{\mathrel{/\mkern-10mu/}}
\DeclareMathOperator{\Maps}{Maps}
\DeclareMathOperator{\Fun}{Fun}
\DeclareMathOperator{\Gr}{Gr}
\DeclareMathOperator{\Pic}{Pic}
\DeclareMathOperator{\Jac}{Jac}
\DeclareMathOperator{\Spec}{Spec}
\DeclareMathOperator{\Vect}{Vect}
\DeclareMathOperator{\Cat}{Cat}
\DeclareMathOperator{\Bord}{Bord}
\DeclareMathOperator{\EndC}{End}
\DeclareMathOperator{\tr}{tr}
\DeclareMathOperator{\sgn}{sgn}
\DeclareMathOperator{\vol}{vol}
\DeclareMathOperator{\Res}{Res}
\DeclareMathOperator{\Tr}{Tr}
\DeclareMathOperator{\Frob}{Frob}
\DeclareMathOperator{\Char}{char}
\DeclareMathOperator{\sdim}{sdim}
\DeclareMathOperator{\rank}{rank}
\DeclareMathOperator{\ord}{ord}
\DeclareMathOperator{\Fns}{Fns}
\DeclareMathOperator{\Hecke}{Hecke}
\DeclareMathOperator{\val}{val}

% ---------- 插图 ----------
% 编号图（原文 Figure n.m）
\newcommand{\insfig}[3][0.8\textwidth]{%
  \begin{figure}[htbp]\centering
  \includegraphics[width=#1]{#2}%
  \caption{#3}\label{fig:#2}%
  \end{figure}}
% 原文未编号的内联示意图：不浮动、不编号，随文走
\newcommand{\inlfig}[2][0.42\textwidth]{%
  \par\vspace{0.45em}\noindent\hfill\includegraphics[width=#1]{#2}\hfill\vspace{0.45em}\par}
% 原文照片
\newcommand{\insphoto}[3][0.75\textwidth]{%
  \begin{figure}[htbp]\centering
  \includegraphics[width=#1]{#2}%
  \caption{#3}\label{fig:#2}%
  \end{figure}}

% ---------- 页眉页脚 ----------
% 中文书刊惯例：正文页码居中置底，不用运行标题（避免 article 的 \MakeUppercase 把
% 西文人名写成全大写，如 Yiannis → YIANNIS）
\pagestyle{plain}
